% P05.1 noncompact literal nonregularity proof, 26 September 2026.

\begin{proof}
The function $\varphi$ is literal piecewise affine and polyhedrally convex.
Positivity at infinity makes $\gamma$ tend to $+\infty$ along every
unbounded face, so $D=(C-\gamma)_+$ is nonnegative, bounded, and compactly
supported.  Refine the original structure by the rational hyperplanes on
which $\gamma=C$, and then take a finite rational simplicial refinement.
We obtain a structure $\Pi$ on which $D$ is affine and for which
\begin{equation}
             D=0\quad\text{on every unbounded cell of }\Pi.
 \label{eq:frontier-polyhedral-noncompact-unbounded-zero}
\end{equation}
Indeed, on any original unbounded cell the recession slope of $\gamma$ is
strictly positive in every nonzero recession direction, so its intersection
with $\{\gamma\leq C\}$ is bounded.
The structure $\Pi$ has finitely many vertices, and every bounded cell is
the convex hull of its vertices.

The choice $C>\gamma(x)$ gives $D(x)>0$, whereas $D$ vanishes near infinity
in the same connected component.  Hence $D$ is nonconstant there.  On some
bounded positive-dimensional cell $\sigma$ its linear part
$k\in M_\sigma$ is therefore nonzero.  Choose a primitive
$w\in N_\sigma$ with $k(w)\neq0$.  Since $\sigma$ is rational, it contains a
rational segment $I=x_0+[0,L]w$ in its relative interior, with $L>0$
rational.  Fix an integer $N>|k(w)|$ and set
\[
                  d:=\frac{D}{N},
                  \qquad P:=\gamma+d.
\]
The function $P$ is polyhedrally convex because it is the displayed convex
combination of $\gamma$ and $\varphi$, and it has $\gamma+O(1)$ growth.

We next construct a bounded literal obstacle $u_N\geq d$ which agrees with
$d$ at every vertex of $\Pi$.  Choose a lattice basis
$m_1,\ldots,m_n$ of $M$ and, for each vertex $v\in\Pi^{(0)}$, define
\[
        \rho_v(y):=\max_{1\leq j\leq n}|m_j(y-v)|.
\]
This is a nonnegative literal piecewise-affine function on $N_{\mathbf R}$,
it vanishes only at $v$, and its affine constants are rational.  For every
vertex $v$ there is an integer $A_v>0$ such that
\begin{equation}
                 a_v:=d(v)+A_v\rho_v\geq d
                 \quad\text{on }\mathbf X.
 \label{eq:frontier-polyhedral-noncompact-cone-majorant}
\end{equation}
To see this, let $K$ be the compact support of $d$.  Near $v$, the finite
piecewise-affine function $d$ is Lipschitz with respect to the norm
$\rho_v$; away from $v$, compactness of $K$ bounds
$(d(y)-d(v))_+/\rho_v(y)$.  Outside $K$ one has $d=0\leq d(v)+A_v\rho_v$.
Increasing $A_v$ to an integer proves
\eqref{eq:frontier-polyhedral-noncompact-cone-majorant}.

Set
\begin{equation}
        u_N:=\min\bigl\{D,\ a_v:\ v\in\Pi^{(0)}\bigr\}.
 \label{eq:frontier-polyhedral-noncompact-obstacle}
\end{equation}
A finite minimum of literal piecewise-affine functions is literal
piecewise affine after a finite rational refinement.  Since $D\geq d$ and
every $a_v\geq d$, while $0\leq D$, the function in
\eqref{eq:frontier-polyhedral-noncompact-obstacle} satisfies
\begin{equation}
        0\leq d\leq u_N\leq D,
        \qquad u_N(v)=d(v)\quad(v\in\Pi^{(0)}).
 \label{eq:frontier-polyhedral-noncompact-vertex-contact}
\end{equation}
It is therefore bounded.  Moreover,
\eqref{eq:frontier-polyhedral-noncompact-unbounded-zero} gives
$u_N=d=0$ on every unbounded cell.

The inequality $u_N\geq d$ shows that $P$ is an admissible competitor under
$\gamma+u_N$.  Conversely, let $h\leq\gamma+u_N$ be any admissible
polyhedrally convex competitor.  At each vertex $v$ of $\Pi$,
\[
                      h(v)\leq\gamma(v)+d(v)=P(v).
\]
If $\tau$ is a bounded cell of $\Pi$, the restriction $h|_\tau$ is an
ordinary convex function.  Every point of $\tau$ is a convex combination of
its vertices, and $P|_\tau$ is affine, so the vertex bounds imply
$h\leq P$ on $\tau$.  If $\tau$ is unbounded, then
$u_N=d=0$ there and the obstacle inequality itself gives
$h\leq\gamma=P$.  Thus $h\leq P$ on $\mathbf X$, proving
\eqref{eq:frontier-polyhedral-noncompact-literal-envelope}.

Let $s\in M_\sigma$ be the linear part of $\gamma$ on $\sigma$.  The slope
of $P$ along $w$ is
\[
                         q=s(w)+\frac{k(w)}{N}.
\]
Both $s(w)$ and $k(w)$ are integers, and
$0<|k(w)|<N$, so $q\notin\mathbf Z$.  Divide $I$ into two subsegments of
equal integral-affine length, with endpoints $x_0,x_1,x_2$.  For any literal
integral-slope polyhedrally convex function $\psi$, the restriction
$\psi|_I$ is convex and piecewise affine with integral slopes: on every
affine subinterval, its integral ambient slope takes the primitive tangent
vector $w$ to an integer.  Its left
derivative $m$ at $x_1$ is therefore an integer.  If
$m\leq\lfloor q\rfloor$, convexity gives
\[
 \psi(x_1)-\psi(x_0)\leq m\frac L2,
 \qquad P(x_1)-P(x_0)=q\frac L2.
\]
If $m\geq\lceil q\rceil$, monotonicity of the one-sided derivatives gives
\[
 \psi(x_2)-\psi(x_1)\geq m\frac L2,
 \qquad P(x_2)-P(x_1)=q\frac L2.
\]
The two cases exhaust the possible values of $m$.  In either case, the
difference of two pointwise errors has absolute value at least
$L\operatorname{dist}(q,\mathbf Z)/2$.  One of those errors has absolute
value at least half that quantity, proving
\eqref{eq:frontier-polyhedral-noncompact-literal-separation}.

The exact envelope is uniformly separated from every literal
piecewise-affine polyhedrally convex function.  It cannot belong to the
literal regularizable class, so the bounded literal obstacle $u_N$ violates
condition~(1) of \cite[Theorem~6.4]{BMPS26}.  Hence literal regularity fails.
Taking $\mathbf X=N_{\mathbf R}$ gives the final assertion.
\end{proof}
